% !TEX program = pdflatex
\documentclass[11pt]{article}
\usepackage[margin=1in]{geometry}
\usepackage{amsmath,amssymb}
\usepackage{enumitem}
\usepackage{booktabs}
\usepackage{xcolor}
\usepackage{titlesec}
\usepackage{parskip}

\definecolor{seccolor}{HTML}{1565C0}
\definecolor{figcolor}{HTML}{2E7D32}
\definecolor{notecolor}{HTML}{B71C1C}

\titleformat{\section}{\large\bfseries\color{seccolor}}{\thesection}{0.6em}{}[\titlerule]
\titleformat{\subsection}{\normalsize\bfseries}{\thesubsection}{0.5em}{}

\newcommand{\fig}[1]{{\color{figcolor}\textsf{[Figure: #1]}}}
\newcommand{\opn}[1]{{\color{notecolor}\textsf{[#1]}}}
\newcommand{\Lb}{\bar\Lambda}

\title{\textbf{Junior Paper Planner}\\[0.3em]
  \large Classification and Measurement of\\
  One-Dimensional Hyperuniform Point Patterns}
\author{Michael Fang}
\date{March 2026}

\begin{document}
\maketitle

%==============================================================================
\section{Thesis}
%==============================================================================

The coarse three-class scheme for hyperuniformity ($\alpha > 1$, $= 1$,
$< 1$) does not distinguish between patterns that differ substantially in
their spatial structure. This paper develops a two-axis framework for
classifying and comparing one-dimensional hyperuniform point patterns:
\begin{itemize}[noitemsep]
  \item $\alpha$ (the hyperuniformity exponent) determines which
        universality class a pattern belongs to;
  \item $\Lb$ (the surface-area coefficient) ranks residual disorder
        \emph{within} each class.
\end{itemize}
We build the first systematic numerical catalog of both metrics across
substitution tilings, perturbed lattices, and stealthy configurations,
computing $\Lb$ for patterns where it was previously unknown and extending
the framework to Classes~II and~III via generalized scaling coefficients.
The catalog reveals what controls each metric: the eigenvalue structure
of the substitution matrix constrains which $\alpha$ values are accessible,
while $\Lb$ is sensitive to the full spatial arrangement of tiles---patterns
sharing the same $\alpha$ can differ in $\Lb$ by more than an order of
magnitude.

%==============================================================================
\section{Paper Structure}
%==============================================================================

%------------------------------------------------------------------------------
\subsection{\S1. Introduction \textnormal{\small (written)}}

\begin{itemize}[leftmargin=1.5em]
  \item Hyperuniformity: $S(k) \to 0$ as $k \to 0$.
        Contrast with Poisson. Where it appears.
  \item Three-class scheme ($\alpha$) is coarse: two Class~I patterns
        can share $\alpha = 3$ but differ in fluctuation amplitude.
  \item $\Lb$ captures this amplitude. Only two prior values known
        (lattice $1/6$, Fibonacci $0.20110$).
  \item What this paper does and main finding: $\alpha$ and $\Lb$
        are independent.
\end{itemize}

\fig{Two Class~I chains with same $\alpha$, different structure}

%------------------------------------------------------------------------------
\subsection{\S2. Theory \textnormal{\small (written)}}

Four subsections:

\begin{itemize}[leftmargin=1.5em]
  \item \textbf{\S2.1. Number variance, structure factor, and $\alpha$.}
        $\sigma^2(R)$, $g_2(r)$, $S(k)$, hyperuniformity condition,
        power-law $S(k) \sim |k|^\alpha$, three classes.

  \item \textbf{\S2.2. The surface-area coefficient $\Lb$.}
        Running average of $\sigma^2(R)$. Lattice minimum $1/6$.
        $\Lambda_{II}$, $\Lambda_{III}$ for Classes~II/III.

  \item \textbf{\S2.3. Substitution tilings and the eigenvalue formula.}
        Substitution rules, matrix $M$, Fibonacci example.
        Eigenvalue formula $\alpha = 1 - 2\ln|\lambda_2|/\ln\lambda_1$.
        Pisot property $\to$ Bragg peaks (dense, countable).
        Cannot fit $S(k)$ directly $\to$ motivates spreadability.

  \item \textbf{\S2.4. Diffusion spreadability.}
        Two-phase embedding, spectral density
        $\tilde\chi_V(k) = \rho|\tilde m(k)|^2 S(k)$.
        Spreadability formula (integral form, $d$ dimensions).
        Gaussian kernel as low-pass filter smooths Bragg peaks.
        Power-law decay $t^{-(d+\alpha)/2}$; physical meaning:
        larger $\alpha$ $\to$ faster diffusion \emph{and} more
        transparent to low-frequency waves.
\end{itemize}

%------------------------------------------------------------------------------
\subsection{\S3. Methods}

Three subsections. Procedural only; all theory/motivation is in \S2.

\begin{itemize}[leftmargin=1.5em]
  \item \textbf{Pattern generation.}
        Substitution tilings: iterative rule application, $\geq 10^6$ tiles.
        URL: $N = 50{,}000$, 17 values of $a \in [0.05, 2.0]$.
        Stealthy: collaborator ensemble, ${\sim}4300$ configs/$\chi$,
        $N = 2000$, $\rho = 1$, $\chi \in \{0.1, 0.2, 0.3\}$.

  \item \textbf{Extraction of $\alpha$ and $\Lb$.}
        $\alpha$: embed in two-phase medium ($\phi_2 = 0.35$),
        compute $S(k)$ via DFT, form $\tilde\chi_V(k)$,
        evaluate spreadability, fit log-log slope.
        Cross-check against eigenvalue formula.
        $\Lb$: sliding window ($30{,}000$ placements, $R < L/4$),
        running average over upper two-thirds of $R$ range.
        Class~II: fit $\sigma^2 = \Lambda_{II}\ln R + b$.
        Class~III: fit $\sigma^2 = A\,R^{\beta}$, both free.

  \item \textbf{Validation.}
        (i) Poisson $\sigma^2 = 2\rho R$, $1.1\%$ error;
        (ii) URL $\Lb = (1+a^2)/6$, $<0.3\%$;
        (iii) Fibonacci $\Lb = 0.200$ vs.\ $0.20110$ (Zachary 2009);
        (iv) Fibonacci $\alpha = 3.05$ via spreadability (Hitin-Bialus 2024).
\end{itemize}

%------------------------------------------------------------------------------
\subsection{\S4. Results}

Single results section with four subsections. Subsections 4.2 and 4.4
carry the most weight.

%- - - - - - - - - - - - - - - - - - - - - - - - - - - - - - - - - - - -
\subsubsection{\S4.1. Metallic-mean family}

Present the main family: all five metallic-mean chains have $\alpha = 3$,
confirmed by both eigenvalue formula and spreadability.

\begin{itemize}[leftmargin=1.5em]
  \item Table of chains with $n$, $\mu_n$, $\rho$, $\alpha$ (eigenvalue),
        $\alpha$ (spreadability), $\Lb$.
        Integer lattice included as baseline ($\Lb = 1/6$).
  \item All are ordered, deterministic quasicrystals with pure
        Bragg spectra. Despite very different tile-length ratios,
        they all share $\alpha = 3$, demonstrating that $\alpha$ is
        determined by the eigenvalue structure alone.
  \item $\Lb$ increases monotonically from Fibonacci ($0.201$) to
        Nickel ($0.310$): same $\alpha$, different residual disorder.
        A first indication that $\alpha$ and $\Lb$ capture
        independent information (confirmed more broadly in \S4.4).
\end{itemize}

\fig{$\sigma^2(R)$ for Fibonacci/Silver/Bronze (bounded Class~I oscillation)}

%- - - - - - - - - - - - - - - - - - - - - - - - - - - - - - - - - - - -
\subsubsection{\S4.2. Structural constraints on $\alpha$}

\textbf{This is a central section of the paper.} What values of $\alpha$
are algebraically accessible, and why?

\begin{itemize}[leftmargin=1.5em]
  \item \textbf{The determinant bound.}
        For a rank-$r$ unimodular Pisot matrix ($|\det M| = 1$):
        $\alpha \leq 1 + 2/(r{-}1)$.
        Rank 2 $\to$ $\alpha \leq 3$ (saturated by all metallic means).
        Rank 3 $\to$ $\alpha \leq 2$.
        Derive from the eigenvalue formula and $|\det M| = 1$.

  \item \textbf{Rank-3 results.}
        Bombieri-Taylor: $\alpha = 1.545$, $\Lb = 0.377$. Occupies
        the interval $1 < \alpha < 2$.
        504 cubic matrices with complex eigenvalue pairs all give
        $\alpha = 2$ exactly (forced by $|\det M| = \lambda_1|\lambda_2|^2 = 1$).
        Example: $a \to bc$, $b \to c$, $c \to abc$; $\Lb = 0.275$.
        Exhaustive enumeration up to rank 4 confirms no unimodular
        matrix gives $\alpha \in (2,3)$.

  \item \textbf{Filling the gap: $|\det M| > 1$.}
        For rank-2: $\alpha = 3 - 2\ln|\det M|/\ln\lambda_1$.
        $|\det M| = 2$ with $\lambda_1 > 4$ gives $\alpha \in (2,3)$.
        138 matrices found (row sums $\leq 10$), 12 distinct $\alpha$
        values in $[2.07, 2.39]$.
        Smallest example: $M = \bigl[\begin{smallmatrix}0&1\\2&4\end{smallmatrix}\bigr]$,
        $\alpha = 2.071$.
        Within each $\alpha$ group, $\Lb$ varies by factors of 5--15,
        showing that the full substitution rule matters, not just the
        eigenvalues.
\end{itemize}

\fig{BT analysis: variance, $S(k)$, spreadability}

\fig{$\sigma^2(R)$ and running $\Lb$ for the $\alpha = 2.071$ chain}

%- - - - - - - - - - - - - - - - - - - - - - - - - - - - - - - - - - - -
\subsubsection{\S4.3. Classes~II and~III}

Brief. Two substitution tilings that fall outside Class~I, included
to complete the picture.

\begin{itemize}[leftmargin=1.5em]
  \item Period-Doubling: $|\lambda_2| = 1$, so $\alpha = 1$ exactly
        (Class~II boundary). $\sigma^2 \sim \Lambda_{II}\ln R$;
        $\Lambda_{II} = 0.080$.
  \item 0222 chain: $|\lambda_2| > 1$ (non-Pisot), $\alpha = 0.639$
        (Class~III). $\sigma^2 \sim \Lambda_{III} R^{1-\alpha}$;
        $\Lambda_{III} = 0.209$.
  \item $\Lambda_{II}$ and $\Lambda_{III}$ serve as the analogues of
        $\Lb$ for these classes: they extract the leading amplitude
        from the appropriate growth law.
\end{itemize}

\fig{$\sigma^2(R)$ for Period-Doubling and 0222 (Class~II/III growth)}

%- - - - - - - - - - - - - - - - - - - - - - - - - - - - - - - - - - - -
\subsubsection{\S4.4. Ranking Class~I patterns by $\Lb$}

\textbf{This is the other central section.} $\alpha$ classifies;
$\Lb$ ranks. The catalog table is presented here, and then we ask:
what do we learn from comparing $\Lb$ across structurally different
systems?

\begin{itemize}[leftmargin=1.5em]
  \item \textbf{The catalog.}
        Full table of all Class~I patterns sorted by $\Lb$:
        integer lattice ($1/6$), metallic means, cubic $\alpha = 2$,
        URL at various $a$, BT, stealthy at various $\chi$.
        Classes~II/III listed separately with $\Lambda_{II}$,
        $\Lambda_{III}$.

  \item \textbf{Comparison families.}
        URL: $\Lb = (1+a^2)/6$, exact (Klatt et al.\ 2020);
        $\alpha = 2$ for all $a > 0$; at $a = 1$ (cloaking),
        $\Lb = 1/3$.
        Stealthy: $S(k) \equiv 0$ for $|k| < K^*$;
        $\Lb \approx 1/(\pi^2\chi)$ for the disordered phase;
        $\chi$ controls the fraction of constrained Fourier modes.

  \item \textbf{What the ranking reveals.}
        The metallic-mean chains are ordered (Bragg spectra,
        $\alpha = 3$); the URL is disordered (continuous $S(k)$,
        $\alpha = 2$). Yet the large-$n$ metallic means approach
        $\Lb \approx 1/3$ from below, matching the cloaked URL.
        Two structurally different systems converge to the same
        residual fluctuation level despite different $\alpha$.

        The cubic $\alpha = 2$ chain ($\Lb = 0.275$) is a
        deterministic analogue of URL at $a \approx 0.8$: same
        exponent, comparable disorder, completely different mechanism.

        Stealthy at $\chi = 0.3$ ($\Lb = 0.357$) has more residual
        disorder than any metallic mean, despite perfect
        small-$k$ suppression. Constraining a finite set of Fourier
        modes does not impose local spatial order.

        The 138 non-unimodular chains show $\Lb$ varying by
        factors of 5--15 within each $\alpha$ group: the full
        substitution rule, not just the eigenvalues, controls
        residual disorder.

        \opn{We only have good ordered structures (metallic means)
        at $\alpha = 3$. A disordered $\alpha = 3$ comparison would
        be valuable but is not in the catalog. Note as limitation.}

  \item \textbf{What controls $\Lb$.}
        Close with a short synthesis: $\Lb$ increases with
        tile-length disparity (metallic means),
        displacement amplitude (URL),
        weaker spectral constraints (stealthy), and
        slower $S(k)$ suppression near $k = 0$ (BT).
        In each case, more room for large-scale density fluctuations
        translates to larger $\Lb$.
\end{itemize}

\fig{Full $\Lb$ catalog table (all patterns, sorted by $\Lb$)}

\fig{$\Lb(\mu_n)$ convergence plot, $n = 1$ through $100$}

%------------------------------------------------------------------------------
\subsection{\S5. Conclusion}

\begin{itemize}[leftmargin=1.5em]
  \item Return to the main thesis: 1D hyperuniformity needs both a
        classification axis ($\alpha$) and a ranking axis ($\Lb$).

  \item Summarize how substitution-matrix algebra constrains $\alpha$:
        determinant condition, rank dependence, $|\det M| > 1$ opens
        new ranges.

  \item Summarize how structural disorder is reflected in $\Lb$:
        tile disparity, displacement amplitude, spectral constraints.

  \item Highlight the physical significance: $\alpha$ controls how
        easily diffusion (and, by extension, wave propagation) occurs
        through the corresponding two-phase medium. Ranking by $\Lb$
        within each class further differentiates patterns that would
        otherwise appear equivalent from a transport standpoint.

  \item Ordered vs.\ disordered: the metallic-mean chains provide a
        clean ordered family at $\alpha = 3$, but the catalog lacks a
        comparably clean disordered family at the same $\alpha$.
        This asymmetry limits how much we can say about the
        ordered/disordered distinction within Class~I.

  \item Open questions:
        \begin{itemize}[noitemsep]
          \item Analytical formula for metallic-mean $\Lb$?
          \item Exact limit of the metallic-mean $\Lb$ sequence?
          \item Extension to 2D (Penrose, Ammann-Beenker)?
          \item What structural features minimize $\Lb$ beyond the lattice?
          \item Comparison with disordered $\alpha = 3$ systems?
        \end{itemize}
\end{itemize}

%==============================================================================
\section{Reference Material}
%==============================================================================

\begin{itemize}[leftmargin=1.5em]
  \item \texttt{jp/jp\_hyperuniformity.tex} --- rough draft of theory/methods;
        has usable equations and definitions to pull from.
  \item \texttt{jp/jp\_outline.md} --- high-level structural outline
        (merged into this planner).
  \item \texttt{results/progress\_report.tex} --- progress report
        with all numerical results and tables.
  \item \texttt{results/figures/} --- all computed figures.
  \item \texttt{results/*.json} --- raw numerical data.
  \item \texttt{literature/} --- all cited PDFs.
\end{itemize}

\end{document}
